\documentclass[12pt]{article}
\usepackage{graphicx}
\usepackage{enumitem}
\usepackage{amsmath}
\usepackage{gvv-book}
\usepackage{gvv}

\title{\textbf{7.4.15}}
\author{\textbf{EE25BTECH11008 - Anirudh M Abhilash}}
\date{October 11, 2025}

\begin{document}

\maketitle

\section*{Question}

The lines $2x - 3y = 5$ and $3x - 4y = 7$ are diameters of a circle having area $154$ sq.units. Find the equation of the circle.

\section*{Solution}

The equations of the diameters are
\begin{align}
2x - 3y &= 5, \\
3x - 4y &= 7.
\end{align}
\begin{align}
\myvec{2 & -3 \\ 3 & -4}\vec{c} &= \myvec{5 \\ 7},
\end{align}
Substituting $\vec{u} = -\vec{c}$:
\begin{align}
\myvec{2 & -3 \\ 3 & -4}\vec{u} &= -\myvec{5 \\ 7}.
\end{align}

Augmented matrix of the system:
\[
\augvec{2}{1}{2 & -3 & -5 \\ 3 & -4 & -7}.
\]

\textbf{Step 1:}
\[
R_2 \to 3R_1 - 2R_2
\]
\[
\augvec{2}{1}{2 & -3 & -5 \\ 0 & -1 & -1}.
\]

\textbf{Step 2:}
\[
R_2 \to -R_2
\]
\[
\augvec{2}{1}{2 & -3 & -5 \\ 0 & 1 & 1}.
\]

\textbf{Step 3:}
\[
R_1 \to R_1 + 3R_2
\]
\[
\augvec{2}{1}{2 & 0 & -2 \\ 0 & 1 & 1}.
\]

\textbf{Step 4:}
\[
R_1 \to \frac{1}{2} R_1
\]
\[
\augvec{2}{1}{1 & 0 & -1 \\ 0 & 1 & 1}.
\]

Thus the solution is
\begin{align}
\vec{u} &= \myvec{-1 \\ 1}.
\end{align}

From the area,
\begin{align}
\pi r^2 &= 154 \quad\Rightarrow\quad r^2 = 49 \quad\Rightarrow\quad r = 7.
\end{align}

Compute \(f\) (using \(\,f=\|\vec{u}\|^2-r^2\)\,):
\begin{align}
\|\vec{u}\|^2 &= (-1)^2 + 1^2 = 2, \\
f &= 2 - 49 = -47.
\end{align}

The general quadratic form is
\begin{align}
\vec{x}^T V \vec{x} + 2\vec{u}^T\vec{x} + f &= 0,
\end{align}
and for a circle \(V = I_{2\times 2}\). Substituting the values:
\begin{align}
\vec{x}^T I \vec{x} + 2\myvec{-1 & 1}\vec{x} - 47 &= 0.
\end{align}
Expanded,
\begin{align}
x^2 + y^2 - 2x + 2y &= 47,
\end{align}
\[
\boxed{x^2 + y^2 - 2x + 2y = 47}
\]

\begin{figure}[H]\centering
\includegraphics[width=1\columnwidth]{figs/plt.png}
\caption{Given Circle}
\label{fig:plt}
\end{figure}

\end{document}
